% !TEX root = ../22830110_dissertation.tex
\chapter{Introduction}
\label{chap:introduction}

\section{Background and Motivation}
Financial reporting forms the cornerstone of corporate accountability. Under frameworks such as the International Financial Reporting Standards (IFRS) and UK GAAP, year-end accounting notes are essential for providing the narrative context required to interpret balance sheets and income statements. However, the preparation of these notes remains a manual, highly labour-intensive process. Accounting professionals routinely consume significant billable hours navigating the meticulous demands of data verification, reconciliation, and compliance alignment. This inefficiency not only escalates costs—making detailed reporting financially unviable for smaller entities—but also heightens the risk of human error, potentially undermining audit integrity.

The emergence of Large Language Models (LLMs) offers a theoretical pathway to automation. Yet, the stochastic nature of these models introduces a critical paradox: accounting requires 100\% numeric fidelity, while LLMs are prone to "hallucinations" and mathematical approximations. Bridging this gap is not merely a task of prompt engineering, but a fundamental challenge of system architecture.

\section{The Research Gap: From Stochastic Text to Deterministic Audit}
While recent advancements in Artificial Intelligence have revolutionized fields from medical imaging to creative writing, its adoption in high-stakes financial auditing has been hampered by a significant research gap. Current AI systems in finance often focus on isolated components—such as anomaly detection or simple data extraction—but fail to integrate deterministic data engineering with constrained generative synthesis. 

Existing Literature often treats LLMs as black-box calculators. There is a documented lack of frameworks that:
\begin{enumerate}
    \item Enforce absolute numeric precision (avoiding floating-point errors) in AI-generated narratives.
    \item Implement multi-layered human-in-the-loop verification that mimics the hierarchical structure of professional firms.
    \item Provide granular, transaction-level provenance for every generated statement.
\end{enumerate}
This dissertation addresses this gap by introducing a hybrid architecture that treats the LLM not as a calculator, but as a \textit{narrative synthesizer} governed by a deterministic rules engine.

\section{Research Question and Objectives}
The central research question of this study is: \textit{To what extent can a multi-agentic AI framework, constrained by deterministic data engineering, automate the production of IFRS/GAAP-compliant accounting notes without compromising numeric accuracy or auditability?}

To address this, the project pursues the following objectives:
\begin{enumerate}
    \item To design and implement a \textbf{Decimal Mandate} architecture for 100\% precise data extraction from the Xero API.
    \item To develop a \textbf{Multi-Agentic Orchestration} layer using CrewAI to simulate professional verification workflows.
    \item To conduct a \textbf{Benchmarking Study} across 10 distinct LLMs to identify the optimal architecture for financial narrative synthesis.
    \item To evaluate the system through a \textbf{Human-in-the-Loop study} with professional accountants, measuring time savings and professional acceptability.
\end{enumerate}

\section{Contribution and Significance}
The primary contribution of this research is the \textbf{ai\_review\_notes} framework, an end-to-end pipeline that demonstrates the viability of "Explainable AI by Design" in professional accounting. Key innovations include:
\begin{itemize}
    \item \textbf{The Decimal Mandate:} A novel data engineering methodology that isolates LLMs from mathematical reasoning, ensuring numeric outputs are sourced only from deterministic code.
    \item \textbf{Agentic Separation of Concerns:} A framework where specialised AI agents (Analyst, Writer, Reviewer) provide the checks and balances necessary for audit-grade outputs.
    \item \textbf{Practical Workflow Evidence:} Survey and benchmark evidence showing that year-end note drafting is time-consuming in practice and that AI-generated drafts are useful when kept inside a professional review loop.
\end{itemize}

\section{Dissertation Structure}
The remainder of this dissertation is structured as follows: Chapter \ref{chap:literature_review} provides a critical review of the state-of-the-art in financial AI and XAI. Chapter \ref{chap:aim_and_objectives} further defines the technical aims. Chapter \ref{chap:analysis_design} details the system's architectural layers. Chapter \ref{chap:methodology} deep-dives into Transformer theory and the comparative test harness for the GPT-5.4 prompting baseline and the live GPT-4.1 fine-tuned deployments. Chapter \ref{chap:implementation} discusses the technical realisation, followed by Chapter \ref{chap:testing_evaluation} which presents the experimental results. Finally, Chapter \ref{chap:conclusions} summarises the findings and outlines future research pathways.
 
